\section{Finite verification and completion of the proof}\label{sec:finite}
\subsection{A common enumeration procedure}\label{sec:enumeration}
Each finite test follows the same sequence: fix an energy and a permitted color partition; assemble the possible majority support components; label their edges from \eqref{eq:entry-catalogue}; discard components using \cref{sec:moments}; and apply the Schur, arithmetic, and two-sided tests to the remaining assemblies. The exact determinant test is always $D>B$, not $D\ge B$.

Here is why the support and phase enumerations are complete. A connected support of order $v$ contains a spanning tree. List all nonisomorphic trees on $v$ vertices, add every subset of missing edges allowed by the energy budget, and retain every resulting graph up to verified isomorphism. Small graphs may instead be read from the complete graph atlas through seven vertices. For each graph, list every allowed norm partition, every placement of those norms on edges, and every associate class of a norm. Fix phases on a spanning tree by diagonal unit switching and list all six residual unit choices on each chord. This switching is used only for spectral and determinant necessary conditions; it is not asserted to be third-root equivalence of the original matrix. When a factorization test uses actual third-root rows, their original row cosets must be retained.

Component records contain their characteristic polynomial, determinant, and the minimum diagonal entry of the zero-eigenvalue projection. For a fixed polynomial, retaining the largest observed minimum projection diagonal is a safe relaxation, not a claim that the components are equivalent. Multisets of component records are assembled with all possible isolate counts. A leaf-containing component is omitted only after \cref{lem:projection} applies, or when the remaining vertices cannot be covered within the energy budget. In the latter case the minimum additional cost is one unit per two vertices.

The resulting computations are summarized below by input family, residual family, and exclusion. Their implementations are identified in \cref{app:replay}. The catalogue-generation arguments are part of the proof; a stored status word is not a substitute for them.

\subsection{Eliminating every size-13 majority block}
\begin{lemma}\label{lem:size13}
No strict improvement has color type $(13,2)$ or $(13,1,1)$ on either side.
\end{lemma}
\begin{proof}
By \cref{thm:orthogonality}, a 13-row same-color block has at least four nonzero internal pairs, hence internal energy at least 36. Its cross energy is at least 78. Therefore $Q\ge114$. For each permitted $Q\ge114$, complete component assembly with \eqref{eq:schur13} leaves no upper bound greater than $B$; these are the size-13 entries of \file{proof/new/exploration.json}, regenerated by \file{proof/new/explore.py}. The component costs satisfy $\sum v_i=13$, $9\sum u_i+c+h=Q$, $c\ge78$, and $h$ comes from the entry catalogue, so this is a finite enumeration with the coverage described above.
\end{proof}
This removes the former energy-0, energy-9, and energy-18 internal cases at once. For example, $P_3$ with ten isolates has 12 independent vertices, and $2K_2$ with nine isolates has 11. The sharp orthogonality theorem excludes them regardless of phases or the two remaining rows. In particular, the 142 size-13 Grams previously subjected to column-equation scans need not be treated one at a time.

After \cref{lem:size13}, a multiple-of-nine shell is pure-color on \emph{both} sides, while a shell congruent to six modulo nine has type $(14,1)$ on both sides. This follows from the cross-pair counts in \cref{lem:colors,lem:partitions}; it does not require choosing an orientation.

\subsection{The mixed-color shells}
The complete one-hub Schur bounds dispose of six shells directly. For reference, \cref{tab:schur} also records the two-hub bounds; all displayed decimals are rounded \emph{up} from exact rational upper bounds.
\begin{table}[ht]
\centering\caption{Upper bounds divided by $B$ from complete component assembly.}\label{tab:schur}
\begin{tabular}{@{}rcc@{}}\toprule
$Q$ & Size-13 majority & Size-14 majority\\\midrule
123&0.965877&0.984600\\
132&0.943735&0.979426\\
141&0.922622&0.971828\\
150&0.896803&0.963221\\
159&0.877308&0.939259\\
168&0.856213&0.919227\\
\bottomrule\end{tabular}
\end{table}

At $Q=87,96,105$, a majority support without isolates needs at least seven edges and hence $Q\ge7\cdot9+14\cdot3=105$. The no-isolate $Q=105$ case has exactly seven disjoint norm-9 edges and minimum cross norms. Equation \eqref{eq:schur14} gives
\[
 \det G\le216^7\left(15-\frac{42}{18}\right)
          =216^7\frac{38}{3}=B.
\]
Thus every strict improvement in these three shells has isolates on both sides. The complete rank-signature enumeration gives the following results.
\begin{table}[ht]
\centering\caption{Residual majority supports after \eqref{eq:mixed-signature}.}\label{tab:mixed}
\begin{tabular}{@{}rll@{}}\toprule
$Q$ & Residual support & Independence number\\\midrule
87 & none & ---\\
96 & $P_3\sqcup C_4\sqcup7K_1$ &11\\
   & $C_4\sqcup10K_1$ &12\\
   & $K_{2,3}\sqcup9K_1$ &12\\
105 & none & ---\\\bottomrule
\end{tabular}
\end{table}
The remaining three supports at $Q=96$ contradict \cref{thm:orthogonality}. The rank catalogue is generated before pairing signatures, using exact characteristic polynomials and kernel-support tests. Its entries retain a lower bound on independence number. The replay checks all compatible row/column signature pairs in both directions. The 52 further Gram candidates retained by the older $Q=96$ phase/cross enumeration are consequently redundant; the new verifier exhibits ten independent vertices in each of them.

For $Q=114$, component bounds leave eight size-14 configurations. Exact core-phase and cross-vector enumeration leaves just one norm-compatible abstract Gram above $B$, of determinant $278010073125000000$. Its majority support is $K_4\sqcup10K_1$. It has an independent set of size 11, so it cannot be a third-root Gram. This also replaces a separate eigenspace-flatness argument. The final cross-vector step examines 3024 states; the complete core screening and the canonical residual Gram are regenerated by \file{proof/new/refine_q114.py}, \file{proof/new/cross_exact.py}, and \file{proof/new/finalize_high.py}.

\subsection{The low pure-color shells \texorpdfstring{$54,63,72,81$}{54,63,72,81}}
For $Q=54$ or 63 there are too few edges to cover all 15 vertices, so both supports have isolates. At $Q=72$, the only no-isolate possibility is $P_3\sqcup6K_2$ with unit normalized edge norms. Its Gram determinant is $3105\cdot216^6$, not an Eisenstein norm. Thus \cref{lem:projection} applies in all three shells, and every active component has minimum degree at least two. Such a component has at most eight edges and vertices. The graph atlas through seven vertices together with $C_8$ is complete in this range. All admissible weight and chord-phase assignments are included.

\begin{table}[ht]
\centering\caption{Exact isolate-case enumeration in the low pure-color shells. The assembly count is after requiring positive kernel projection on every active coordinate.}\label{tab:low}
\begin{tabular}{@{}rrrrr@{}}\toprule
$Q$ & Phase assignments & Assemblies & Norm-compatible $D>B$ & After projection\\\midrule
54&372&7&1&0\\
63&1086&15&1&0\\
72&7524&58&1&0\\
81&42720&285&22&0\\\bottomrule
\end{tabular}
\end{table}
At $Q=81$, a no-isolate support has nine unit-norm edges. It is either a forest with six components, each of order at most five, or $K_3\sqcup6K_2$. The forest cases fail the determinant and norm tests. The four residual triangle spectra have no zero eigenvalue after subtracting $15I$, so a no-isolate Gram cannot be paired with an isolate Gram. For the isolate case, the catalogue extends through nine energy units. The additional eight-vertex, nine-edge minimum-degree-two graphs are the figure-eight, barbell, and theta graphs; $C_8$ and $C_9$ are also included. The exact counts are in \cref{tab:low}.

The no-isolate triangle cases are excluded in \cref{app:q81}. That argument uses the same intertwining identity as the other shells; its last step is an elementary obstruction to a sum of two scaled permutation blocks. It is included in full because a spectral catalogue alone does not exclude the negative-flux triangle.

\subsection{The pure-color shells \texorpdfstring{$90,99,108$}{90,99,108}}
The complete component catalogues extend to ten energy units for $Q=90$ and to twelve units for $Q=99,108$. Large supports are generated by the spanning-tree procedure in \cref{sec:enumeration}, with the no-isolate covering budget and the isolate/leaf restriction applied separately. Table~\ref{tab:pure} records the assembly counts at the last stage; these are not counts of inequivalent original matrices.
\begin{table}[ht]
\centering\caption{Pure-color residual computations.}\label{tab:pure}
\begin{tabular}{@{}rccp{66mm}@{}}\toprule
$Q$ & Isolate assemblies & No-isolate assemblies & Remaining test\\\midrule
90 &807&74&Two no-isolate spectra; spectral rectangle.\\
99 &435&632&14 component--spectrum assemblies, 16 equal-spectrum pairs; spectral rectangle.\\
108&146&5570&No no-isolate norm survivor; eight isolate norm survivors all fail projection.\\\bottomrule
\end{tabular}
\end{table}

At $Q=90$, each surviving no-isolate spectrum consists of five $K_2$ components and one component of order five. Take the five-vertex component on one side and the ten pair coordinates on the other. The rectangle has squared Frobenius norm 50. Its common characteristic factor is respectively $x-1$ or $x^2-1$, giving the impossible bounds $50\le18$ or $50\le30$ by \cref{thm:rectangle}.

At $Q=99$, all isolate cases fail the two-sided projection test. The 14 no-isolate component--spectrum assemblies give 16 unordered pairs with equal total characteristic polynomial. For each pair, all nonempty component unions are examined. Every pair violates \eqref{eq:rectangle}. In particular, the three pairs that survive a weaker rank-times-largest-eigenvalue test give
\[
 40\le30,\qquad20\le18,\qquad36\le30.
\]
The certificate stores the two selected unions, their polynomials, their gcd, and the two integers in the failed inequality. These witnesses are independently re-evaluated by the revision verifier.

At $Q=108$, neither an isolate nor a no-isolate candidate survives. The largest catalogue used for this step contains 15,899,506 normalized phase states. It is regenerated, not accepted from its reported size. The former two-hub energy-18 branch is already removed by \cref{lem:size13}.

\subsection{The remaining pure-color shells}
For $Q=135,144,153,162$, complete component bounds followed by the two-sided projection iteration leave no configuration. At $Q=117$ the iteration leaves
\[
 6K_1+\text{a 9-vertex core},\quad
 7K_1+\text{an 8-vertex core},\quad
 8K_1+\text{a 7-vertex core},
\]
all with core energy 13 units. The respective rank bounds for $T$ are 5, 4, and 4. At $Q=126$ only seven isolates and an 8-vertex core of energy 14 units remain, with rank at most four. These assertions come from component-cost enumeration and the nullity lower bound in \cref{lem:projection}, not from an assumption about the equality structure.

The characteristic polynomial of an integral Hermitian Eisenstein matrix has integer coefficients: its coefficients lie in $\Eis$ and are real. If $\rank T\le d$, it follows that
\begin{equation}\label{eq:lattice}
 \det G=15^{15-d}3^dR,\qquad R\in\mathbb Z.
\end{equation}
Extra zero roots may be included when the actual rank is smaller. At $Q=126$, the rank-four envelope and $D>B$ give $397\le R\le401$. Every integer $15^{11}3^4R$ in this interval fails the Eisenstein norm test. The test is on the \emph{full determinant}; it is not a claim that none of the integers $R$ is a norm.

At $Q=117$, $\tr T^2=26$. In the rank-four case the only norm-compatible determinant quotient is $R=405$. The nonzero factor, padded by zero roots when necessary, has the form
\[
 x^4-13x^2-e_3x+e_4,\qquad
 |e_3|\le22,\qquad e_4=405-300-5e_3.
\]
The triangle bound gives the stated bound on $e_3$. None of these 45 polynomials has four real roots. In the rank-five case, the moment refinement leaves $R\in\{1983,1984,1987\}$, and the factor is
\[
 x^5-13x^3-e_3x^2+e_4x-e_5,
\]
where
\[
 |e_3|\le12,\qquad -85\le e_4\le50,\qquad |e_5|\le62,
 \qquad R=1500+25e_3+5e_4+e_5.
\]
Here the triangle bound is six, so $|e_3|\le12$. Newton identities give $e_4=(26^2-2\tr T^4)/8$, and the second moment bounds its range. The product bound $|e_5|\le(26/5)^{5/2}<62$ gives the last coefficient bound. Exactly 240 coefficient tuples meet the displayed integer restrictions; none has five real roots. The verifier uses rational real-root isolation, so failure of real-rootedness is an exact contradiction to Hermitian symmetry.

\subsection{Assembly of the global argument}
\begin{proof}[Proof of \cref{thm:main}]
Suppose $\det(HH^*)>B$. By \cref{cor:shells}, its energy lies in the 23-element set \eqref{eq:qset}. Lemmas~\ref{lem:partitions} and \ref{lem:size13} reduce the color types to pure color on both sides or $(14,1)$ on both sides. The mixed-color shells are excluded by \cref{tab:schur,tab:mixed} and the $Q=114$ independent-set obstruction. The pure-color shells are excluded by the low-energy catalogues, \cref{app:q81}, the spectral rectangles, and the rank/moment computations above. These cases partition \eqref{eq:qset}. Thus no strict improvement exists. The explicit matrix in \cref{app:matrix} has squared determinant $B$, as verified independently in \cref{sec:witness}, proving equality.
\end{proof}
